\chapter{89}

\section[Mailand (IT), Sonntag, 9. September]{Mailand (IT), Sonntag, 9. September}

\subsection{Valentina}

Wo hört Loyalität auf, und wo beginnt meine Sorge um mich selbst?

Die Frage macht mich wahnsinnig. Maria und Franco sind sich einig. Ich
musste es ihnen erzählen. Die Geschichte ist mir endgültig zu schwer
geworden. Ich hänge in einer Angst fest, die mich bald umbringt. Es sei
an der Zeit, mich mit der Polizei zusammenzutun, sagen sie.
Eindringlich. Maria bot mir sogar an, für eine Zeit in ihrer Wohnung
abzutauchen.

Mit welcher Arroganz diese Dreckskerle aufgetaucht sind, plötzlich und
so unwirklich, dass ich immer noch an meiner Wahrnehmung zweifle und
schon beim leisesten Gedanken daran, in Schweiss ausbreche.

Ich schlage die Bettdecke zurück, setze mich auf und greife mir mein
Handy. Seit gestern habe ich es ganz ausgestellt. Meine Finger versteifen sich
und mir graut davor, es zu aktivieren. Ich halte den Atem an und tue es.
Die vielen neuen Nachrichten springen mich an. Auch Maria und Mama
schrieben und Louis wollte mich anrufen. Louis. Schon alleine sein Name
löst ein wohliges Körpergefühl aus. Ähnlich, wie ich es nach einem
Vollbad kenne. Ich muss ihm die Wahrheit sagen. Erzählen, was sich
zugetragen hat. Egal, wie er reagiert. Mir kommen die Tränen. Meine
Angst während des Überfalls wird mit dem weggestopften Gefühl
unerwiderter Zuneigung zu einem explosiven Gemisch. Mit der freien Hand
halte ich mich an der Matratze fest. Ich wähle Louis’ Nummer.

Dass er schon nach dem zweiten Klingeln abnimmt, lässt mich zusammenzucken.

«Valentina?»
